\documentclass[10pt,a4paper]{amsart}
\usepackage[margin=19mm]{geometry}
\usepackage{amsmath,amssymb,mathtools}
\usepackage[scr=rsfs,cal=euler]{mathalfa}
\usepackage{xfrac}
\usepackage[unicode,hidelinks,bookmarks=false]{hyperref}
\newcommand{\ie}{i.e.\ }
\newcommand{\cf}{cf.\ }
\newcommand{\resp}{respectively\ }
\newcommand{\Subsection}{\subsection}
\providecommand{\nobreakdash}{\nobreak\hbox{-}}
\setlength{\emergencystretch}{2em}
\multlinegap0pt
\usepackage{bm}
\usepackage{bbm}
\usepackage{dsfont}
\usepackage{xspace}
\usepackage{cancel}
\usepackage{mathtools}
\usepackage{amsthm}
\makeatletter
\@ifundefined{lemma}{\newtheorem{lemma}{Lemma}}{}
\makeatother
\title[Moment Martingale Posteriors for Semiparametric Predictive B]{Moment Martingale Posteriors for Semiparametric Predictive Bayes}
\author{Yiu Yin Yung, Stephen M.S. Lee, Edwin Fong}
\date{Source: arXiv:2507.18148v1 (CC BY 4.0)}
\begin{document}
\maketitle

\noindent\emph{Source and licence.} Moment Martingale Posteriors for Semiparametric Predictive Bayes. Version arXiv:2507.18148v1 (CC BY 4.0), distributed under the Creative Commons Attribution 4.0 International licence, as recorded in the arXiv licence metadata for this exact version and shown on the abstract page https://arxiv.org/abs/2507.18148v1. Full rights notice: licenses/arxiv-2507.18148-CC-BY-4.0.txt.\par\smallskip

\noindent\textit{Editorial scope.} Counted unit: the Lemma whose source label is \texttt{lem:psi3} in the article (source0.tex lines 1394--1397, source label \texttt{lem:psi3}), delivered together with the article's own complete original proof of it (source0.tex lines 1398--1446, one \texttt{proof} environment of 49 lines). No mathematics is written, repaired or re-derived: statement and proof are the article's own text line for line. Required context retained uncounted: none. Every definition, display and label of those lines is retained unchanged. Omitted from this package and not redistributed: 1--1393, 1447--1905. Nothing inside a retained excerpt is omitted or altered: every retained body is delivered exactly as the article's lines are written, so no omission receipt of any kind accompanies this package. Citations: the retained excerpts carry no citation command at all, so no bibliography entry of the article is transcribed and no reference is redistributed. Reference adaptation: none was needed and none was made; every reference inside the delivered unit resolves inside it, so no unresolved reference or citation is produced. Printed slips kept literal: no printed word, symbol, hypothesis or number of the counted statement or of its proof was altered. Layout and class: the article's own class is \texttt{article} with packages including amsmath, amssymb, fontenc, inputenc, textcomp, lmodern, xcolor, longtable, booktabs, array, calc, etoolbox, footnote, graphicx; the wrapper is the lane's pinned \texttt{amsart} editorial wrapper at 10pt on A4, which restates the theorem-like environments the retained text needs and adds the article's own macro definitions that the retained text needs (none), with extra wrapper packages (bm, bbm, dsfont, xspace, cancel, mathtools). The delivered numbering is the unit's own. Assets: no retained span contains a figure environment, a graphics-inclusion command, a PDF-page inclusion, a table or a TikZ picture; no image file is copied into this package, the package declares no asset dependency and no image file is redistributed with it. Licence: version arXiv:2507.18148v1 of this article is distributed under the Creative Commons Attribution 4.0 International licence, as recorded in the arXiv licence metadata for this exact version and shown on the abstract page https://arxiv.org/abs/2507.18148v1. A full rights notice is delivered at licenses/arxiv-2507.18148-CC-BY-4.0.txt. Characters: every retained excerpt is pure ASCII, so the delivered engine, XeLaTeX with its default Latin Modern text font, reports no missing character.
\begin{lemma}
Assume that there exists a neighborhood $U$ of $\theta^*$ such that $\sup_{\theta \in U}\mu_\theta^{(2)} < \infty$. Then we have
$\Psi_{C,n}\overset{p}{\to} C.$ \label{lem:psi3}
\end{lemma}

\begin{proof}
First, we note that $C \leq 2E_{f_{\theta^*}}[|X|]<\infty$ from the bounded variance assumption. Next, if we can show that for any deterministic sequence $\theta_n$ converging to $\theta^*$, it follows that $\Psi_{C,n}(\theta_n)\to\Psi_{C,n}(\theta^*)=C$, then $\Psi_{C,n}(\theta)$ is continuous at the point $\theta^*$. We can thus apply continuous mapping theorem to show that $\Psi_{C,n}\overset{p}\to C$. 

Now we consider a sequence $\theta_n\to\theta^*$ deterministically to show continuity of $\Psi_{c,n}(\theta)$. Consider a sequence $X_n \sim f_{\theta_n}$ for $n \geq 1$ and $X \sim f_{\theta}^*$. 
We begin by showing a few intermediate results that will be useful in the proof. Note that the assumption implies that $\{X_n\}_{n \geq 1}$ is uniformly integrable as $X_n$ is bounded in $L_2$, and $E_{f_{\theta}^*}[|X|]<\infty$. As a result, for any $\epsilon > 0$, we can choose $K$ sufficiently large so that
\begin{align}
    \int_{|x|> K}|x|\,|f_{\theta_n}(x)-f_{\theta^*}(x)|dx\leq \int_{|x|> K}|x|\,f_{\theta_n}(x)dx+\int_{|x|> K}|x|\,f_{\theta^*}(x)dx <\frac{\epsilon}{2}\cdot \label{eq:ui1}
\end{align}
Next, choose $n$ sufficiently large so that
\begin{align}
    \left\|f_{\theta_n}-f_{\theta^*}\right\|_{L_1} \leq \frac{\epsilon}{2K} \label{eq:L1}
\end{align}
which is possible as $f_{\theta_n}(x) \to f_{\theta^*}(x)$ for each $x$ (from continuity), and Scheff\'{e}'s lemma gives
$$
\|f_{\theta_n}- f_{\theta^*}\|_{L_1} \to 0.
$$
Now we prove that $\Psi_{C,n}\to C$ deterministically by considering the absolute difference 
$$|\Psi_{C,n} -C|=\Big|\iint|x-x'|\left[f_{\theta_n}(x)f_{\theta_n}(x')-f_{\theta^*}(x)f_{\theta^*}(x')\right]dx\,dx'\Big|.$$
We rewrite the term 
$$f_{\theta_n}(x)f_{\theta_n}(x')-f_{\theta^*}(x)f_{\theta^*}(x')=f_{\theta_n}(x)\left[f_{\theta_n}(x')-f_{\theta^*}(x')\right]+f_{\theta^*}(x')\left[f_{\theta_n}(x)-f_{\theta^*}(x)\right].$$
Thus, 
\begin{align*}
|\Psi_{C,n}-C|&\leq \underbrace{\iint|x-x'|f_{\theta_n}(x)\left|f_{\theta_n}(x')-f_{\theta^*}(x')\right|dx\,dx'}_{T_1} \\
&+ \underbrace{\iint|x-x'|f_{\theta^*}(x')\left|f_{\theta_n}(x)-f_{\theta^*}(x)\right|dx\,dx'}_{T_2}.
\end{align*}
For $T_1$, we have 
$$T_1=\int\left(\int|x-x'|f_{\theta_n}(x)dx\right)\left|f_{\theta_n}(x')-f_{\theta^*}(x')\right|dx'.$$
The inner integral is $E_{f_{\theta_n}}[|X_n-x'|]$. Note that $|X_n-x'|\leq |X_n|+|x'|$, so $E_{f_{\theta_n}}[|X_n-x'|]\leq E_{f_{\theta_n}}[|X_n|]+|x'|$. Thus, we have 
\begin{flalign*}T_1&\leq \int\left(E_{f_{\theta_n}}[|X_n|]+|x'|\right)\left|f_{\theta_n}(x')-f_{\theta^*}(x')\right|dx'\\
&=E_{f_{\theta_n}}[|X_n|]\cdot \|f_{\theta_n}-f_{\theta^*}\|_{L_1} + \int|x'|\left|f_{\theta_n}(x')-f_{\theta^*}(x')\right|dx'.
\end{flalign*}
The first term converges to 0 as $\sup_{n} E_{f_{\theta_n}}|X_n| < \infty$ from uniform integrability, and $\|f_{\theta_n}-f_{\theta^*}\|_{L_1} \to 0$ from (\ref{eq:L1}). We will handle the second term shortly.

Similarly, we have 
\begin{flalign*}T_2&\leq E_{f_{\theta^*}}|X|\cdot\|f_{\theta_n}-f_{\theta^*}\|_{L_1} + \int|x|\left|f_{\theta_n}(x)-f_{\theta^*}(x)\right|dx.
\end{flalign*}
Again the first term converges to 0.

Now consider the common second term in $T_1$ and $T_2$, which we can decompose as
$$\int_{|x|\leq K}|x|\,|f_{\theta_n}(x)-f_{\theta^*}(x)|dx+\int_{|x|> K}|x|\,|f_{\theta_n}(x)-f_{\theta^*}(x)|dx.$$
For the first integral we have 
$$\int_{|x|\leq K}|x|\,|f_{\theta_n}(x)-f_{\theta^*}(x)|dx\leq K \int_{|x|\leq K}\,|f_{\theta_n}(x)-f_{\theta^*}(x)|dx \overset{\textnormal{(\ref{eq:L1})}}\leq K  \frac{\epsilon}{2K}=\frac{\epsilon}{2},$$
while for the second integral it is bounded by $\epsilon/2$ by (\ref{eq:ui1}), this means there exist a sufficiently large $n$ such that 
$$\int|x|\,|f_{\theta_n}(x)-f_{\theta^*}(x)|dx \leq \epsilon.$$

As a result, both $T_1$ and $T_2$ converge to 0, which is sufficient to imply $\Psi_{C,n}\to C$ for any deterministic sequence of $\theta_n\to \theta^*$. 


\end{proof}

\end{document}
